\documentclass[a4paper]{article}
\usepackage{titlesec}
\usepackage{geometry}
\geometry{hmargin={2cm,2cm}, vmargin={3.5cm,3.5cm}}
\usepackage{graphicx}
\usepackage[table]{xcolor}
\usepackage{hyperref}
\hypersetup{pdfborder = {0 0 0},
colorlinks,
citecolor=black,
filecolor=black,
linkcolor=black,
urlcolor=blue}
\definecolor{logo}{RGB}{20,99,163}
\usepackage[T1]{fontenc}
\usepackage[utf8]{inputenc}
\usepackage[english,italian]{babel}
\usepackage{fancyhdr}
\usepackage{lastpage}
\usepackage{comment}

% creazione comando subsubsubsection
\titleclass{\subsubsubsection}{straight}[\subsection]

\newcounter{subsubsubsection}[subsubsection]
\renewcommand\thesubsubsubsection{\thesubsubsection.\arabic{subsubsubsection}}
\titleformat{\subsubsubsection}
    {\normalfont\normalsize\bfseries}{\thesubsubsubsection}{1em}{}
    \titlespacing*{\subsubsubsection}
    {0pt}{3.25ex plus 1ex minus .2ex}{1.5ex plus .2ex}

\makeatletter
\def\toclevel@subsubsubsection{4}
\def\l@subsubsubsection{\@dottedtocline{4}{7em}{4em}}
\makeatother

\setcounter{secnumdepth}{4}
\setcounter{tocdepth}{4}
%fine creazione comando


%Comandi di impaginazione uguale per tutti i documenti
\pagestyle{fancy}
\lhead{\includegraphics[scale=0.2]{../../../../Images/logo.png}}
%Titolo del documento
\rhead{\large Verbale Interno del 2021-03-13}
%\rfoot{\thepage}
\cfoot{Pagina \thepage{} di \pageref{LastPage}}
\setlength{\headheight}{35pt}
\renewcommand{\footrulewidth}{0.4pt}
\begin{document}
        
    \begin{titlepage}
        \begin{center}
            \includegraphics{../../../../Images/logo.png}\\
            \vspace{20px}
            \textcolor{logo}{\hrulefill}\\
            \vspace{20px}
            \textbf{\huge\textcolor{logo}{Verbale Interno del 2021-03-13}}\\
            \vspace{10px}
            \textcolor{logo}{\hrulefill}\\
            \vspace{40px}
            \textbf{\Large Informazioni sul documento}\\
            \vspace{20px}
            \begin{tabular}{p{100px} | p{100px}}
                \textbf{Versione} & 0.1\\
                \textbf{Responsabile} & Manuel Varo\\
                \textbf{Redattore} & Nicolò Giaccone\\
                \textbf{Verificatore} & Samuele De Simone\\
                \textbf{Destinatari} & Prof. Tullio Vardanega \newline Prof. Riccardo Cardin, \newline TechSWEave\\
                \textbf{Stato} & Da approvare\\
                \textbf{Uso} & Interno\\
            \end{tabular}\\
            \vspace{60px}
            \href{mailto:techsweave@gmail.com}{techsweave@gmail.com}\\
    
        \end{center}
        \end{titlepage}
    
        %DIARIO DELLE MODIFICHE
        \begin{center}
            \textbf{\Large Diario delle modifiche}\\
            \vspace{10px}
            \begin{table}[h!]
            \centering
            \rowcolors{2}{logo!10}{logo!40}
            \renewcommand{\arraystretch}{1.8}
            \begin{tabular}{p{150px} p{90px} p{70px} p{60px} p{45px}}
                \rowcolor{logo!70} \textbf{Modifica} & \textbf{Autore} & \textbf{Ruolo} & \textbf{Data} & \textbf{Versione}\\
                Stesura documento & Nicolò Giaccone & Analista & 2021-03-14 & 0.0.1\\   
                Verifica & Samuele De Simone & Verificatore & 2021-03-15 & 0.1.0\\
                Validazione & Manuel Varo & Responsabile & 2021-03-16 & 1.0.0
            \end{tabular}
        \end{table}
        \end{center}
    
        %INDICE
        \newpage
        \tableofcontents
        \newpage

    \section{Informazioni riunione}
    \begin{itemize}
        \item Data e ora: 2021-03-13 ore 09.30-10.45.
        \item Luogo: Canale Discord del gruppo tramite videochiamata.
        \item Partecipanti interni: Samuele De Simone, Manuel Varo, Tito Scutari, Nicolò Giaccone, Amedeo Meggiolaro, Simone Urbani, Marco Barbaresco.
        \item Partecipanti esterni: nessuno.
        \item Segretario: Amedeo Meggiolaro.
    \end{itemize}
    \section{Argomenti della riunione}
    \begin{itemize}
        \item Scelta dell'approccio nell'uso di GitHub;
        \item Discussione sugli strumenti da usare per un lavoro agile;
        \item Conoscenza dei documenti da produrre;
        \item Invio mail ai proponenti per fissare il primo incontro;
        \item Divisione gruppi per redazione documenti;
    \end{itemize}
    \section{Resoconto}
    
        \subsection{Scelta dell'approccio all'uso di GitHub}
        In concordo con i partecipanti del gruppo si è scelto di seguire l'approccio \textit{GitHub Flow}, per andare a gestire al meglio il lavoro di gruppo senza andare a creare troppi conflitti tra gli utilizzatori della repository.
        
        \subsection{Discussione sugli strumenti da usare per un lavoro agile}
        Il gruppo ha discusso sul come approcciarsi allo sviluppo del progetto per massimizzare la produttività. Simone Urbani ha prosto l'utilizzo del software \textit{Asana}, per potersi dividere i compiti fissando degli obbiettivi e delle timeline. Amedeo Meggiolario ha proposto di raggiungere lo stesso obbiettivo usando le teconologie messe a disposizione da GitHub, quali \textit{milestone} e \textit{issue}

        \subsection{Conoscenza dei documenti da produrre}
        Il gruppo ha discusso su quali documenti lavorare 
        \begin{itemize}
            \item Norme di progetto;
            \item Analisi dei requisiti;
            \item Piano di qualifica;
            \item Glossario;
        \end{itemize}

        \subsection{Invio mail ai proponenti per fissare il primo incontro}
        Abbiamo inviato una mail ai proponenti per fissare il primo incontro di conoscenza;

        \subsection{Divisione gruppi per redazione documenti}
        Il gruppo ha diviso i vari membri in sotto gruppi per redigere al meglio i documenti da produrre in vista della RR.
\end{document}
